\documentclass[11pt]{article}
\usepackage[a4paper,margin=1in]{geometry}
\usepackage{amsmath,amssymb,amsthm,mathtools}
\usepackage{graphicx}
\usepackage{xcolor}
\usepackage{booktabs}
\usepackage{enumitem}
\usepackage[hidelinks]{hyperref}
\usepackage{microtype}

\definecolor{navy}{RGB}{24,55,92}
\definecolor{darkgreen}{RGB}{22,100,72}
\definecolor{softgray}{RGB}{245,247,250}
\definecolor{softred}{RGB}{252,240,240}

\newtheorem{theorem}{Theorem}
\newtheorem{corollary}[theorem]{Corollary}
\newtheorem{proposition}[theorem]{Proposition}
\newtheorem{lemma}[theorem]{Lemma}
\theoremstyle{definition}
\newtheorem{remark}[theorem]{Remark}

\newcommand{\T}{\mathbb T}
\newcommand{\N}{\mathbb N_0}
\newcommand{\K}{\mathcal K}
\newcommand{\B}{\mathcal B}
\newcommand{\Real}{\operatorname{Re}}
\newcommand{\essinf}{\operatorname*{ess\,inf}}

\title{\textbf{Coercivity of Hermitian and Toeplitz Form Matrices}\\
\large A full positive answer to the structured-matrix question in arXiv:0807.2362}
\author{Automated Mathematics Run \texttt{fa\_banach\_001}\\
Lane 11, model GPT5.6}
\date{9 August 2026}

\begin{document}
\maketitle

\begin{center}
\fcolorbox{darkgreen}{softgray}{\parbox{0.89\textwidth}{
\textbf{Status: candidate full solution, likely valid.}
For bounded Hermitian block forms, scalar Schur-complement arguments transfer
verbatim after Riesz representation.  For bilateral and one-sided Toeplitz
form matrices with absolutely summable operator coefficients, coercivity is
equivalent to uniform positivity of the Hermitian part of the Fourier symbol;
the theorem also gives the exact optimal coercivity constant.
}}
\end{center}

\section{Source question}

The source is Stefano Cardanobile, \emph{Diffusion systems and heat equations
on networks} \cite{Cardanobile}, arXiv:0807.2362 (2008).  In the discussion of
Section 1.3, printed page 49 (PDF page 60), it asks:

\begin{quote}
``An open question is whether it is possible to translate other arguments
holding for scalar-valued matrices in the context of forms, e.g. symmetric
matrices or Toeplitz matrices.''
\end{quote}

\begin{figure}[ht]
  \centering
  \includegraphics[width=\textwidth]{figures/open_problem_crop.png}
  \caption{The open question in the original source, printed page 49.  The
  full readable page width is retained.}
\end{figure}

The wording is existential and deliberately broad.  The results below answer
both named examples: an exact Schur criterion for Hermitian form matrices and
an exact symbol criterion for a natural infinite-dimensional Toeplitz class.

\section{Definitions and scope}

All Hilbert spaces are complex and inner products are linear in the first
variable.  A bounded sesquilinear form $a$ on a Hilbert space $V$ has a unique
Riesz representative $A\in\B(V)$ such that
\[
  a(u,v)=\langle Au,v\rangle.
\]
Its optimal real lower bound is
\[
  c(a):=\inf_{u\ne0}\frac{\Real a(u,u)}{\|u\|^2}.
\]
Thus $a$ is coercive exactly when $c(a)>0$.  The form is Hermitian when
$a(u,v)=\overline{a(v,u)}$, equivalently $A=A^*$.

For the Toeplitz result, let $\K$ be an arbitrary complex Hilbert space and
let $(B_k)_{k\in\mathbb Z}\subset\B(\K)$ satisfy
\begin{equation}\label{eq:wiener}
  \sum_{k\in\mathbb Z}\|B_k\|<\infty.
\end{equation}
This includes every finite-band operator-valued Toeplitz matrix.  For
$J=\mathbb Z$ or $J=\N$, first define on finitely supported vectors
$u,v\in\ell^2(J;\K)$
\begin{equation}\label{eq:toeplitzform}
  a_J(u,v):=\sum_{m,n\in J}
  \langle B_{m-n}u_n,v_m\rangle,
\end{equation}
and then extend continuously.  Condition \eqref{eq:wiener} guarantees this
extension.  Its norm-continuous symbol and Hermitian symbol are
\begin{equation}\label{eq:symbols}
  F(\theta):=\sum_{k\in\mathbb Z}B_ke^{ik\theta},\qquad
  G(\theta):=\frac{F(\theta)+F(\theta)^*}{2}.
\end{equation}

\section{Proof intuition}

For a Hermitian block form, Riesz representation changes the entries from
scalar-valued sesquilinear maps into bounded operators.  Gaussian elimination
then becomes a bounded triangular congruence.  The same multiplication that
produces a scalar Schur complement produces an operator Schur complement, and
bounded invertibility of the triangular factor makes strict positivity
equivalent on both sides.

For Toeplitz forms, the repeated diagonal structure is convolution.  Fourier
series turn convolution into pointwise multiplication by $F(\theta)$, so the
real quadratic form only sees $G(\theta)$.  Uniform positivity of $G$ is
immediately sufficient.  The only delicate direction is necessity for the
one-sided matrix: normalized geometric sequences in $\ell^2(\N;\K)$ have
Fourier mass equal to a Poisson kernel.  Concentrating that mass at an
arbitrary angle and in an arbitrary coefficient-space direction recovers
every pointwise Rayleigh quotient of $G$.  No boundary loss occurs.

\section{Symmetric matrices: the exact Schur criterion}

\begin{theorem}[Hermitian two-block form criterion]\label{thm:schur}
Let $V_1,V_2$ be Hilbert spaces and let a bounded Hermitian form on
$V_1\oplus V_2$ have Riesz representative
\[
 A=\begin{pmatrix}P&B^*\\ B&Q\end{pmatrix},
\]
where $P=P^*$ on $V_1$, $Q=Q^*$ on $V_2$, and
$B\in\B(V_1,V_2)$.  Suppose $P$ is coercive.  Define
\[
 S:=Q-BP^{-1}B^*.
\]
Then the full form is coercive if and only if $S$ is coercive.
\end{theorem}

\begin{proof}
Since $P$ is coercive, $P^{-1}$ is bounded.  Direct block multiplication gives
the exact factorization
\begin{equation}\label{eq:ldl}
 A=
 \underbrace{\begin{pmatrix}I&0\\ BP^{-1}&I\end{pmatrix}}_{L}
 \begin{pmatrix}P&0\\0&S\end{pmatrix}
 \underbrace{\begin{pmatrix}I&P^{-1}B^*\\0&I\end{pmatrix}}_{L^*}.
\end{equation}
The triangular operator $L$ is boundedly invertible.  If
$P\ge pI$ and $S\ge sI$ with $p,s>0$, then for every $z\in V_1\oplus V_2$,
\[
 \langle Az,z\rangle
 \ge \min(p,s)\|L^*z\|^2
 \ge \frac{\min(p,s)}{\|L^{-1}\|^2}\|z\|^2.
\]
Thus $A$ is coercive.

Conversely, suppose $A\ge\gamma I$ with $\gamma>0$.  Restricting to
$V_1\oplus\{0\}$ gives $P\ge\gamma I$.  For $y\in V_2$, insert
$z=(-P^{-1}B^*y,y)$ in the quadratic form.  Formula \eqref{eq:ldl} yields
\[
 \langle Sy,y\rangle=\langle Az,z\rangle
 \ge\gamma\bigl(\|P^{-1}B^*y\|^2+\|y\|^2\bigr)
 \ge\gamma\|y\|^2.
\]
Hence $S$ is coercive.
\end{proof}

\begin{corollary}[Finite Hermitian block elimination]\label{cor:recursive}
For a bounded Hermitian $n\times n$ form matrix on
$V_1\oplus\cdots\oplus V_n$, recursively take the first diagonal pivot and
replace the remaining block by its Schur complement.  The full form is
coercive if and only if every pivot produced by this recursion is coercive.
\end{corollary}

\begin{proof}
Apply Theorem \ref{thm:schur} to the splitting
$V_1\oplus(V_2\oplus\cdots\oplus V_n)$ and iterate on the Schur complement.
If the full form is coercive, each leading compression and each resulting
Schur complement is coercive by the converse part of the theorem.  If every
pivot is coercive, reverse the bounded triangular factorizations.
\end{proof}

This is precisely the operator/form version of Hermitian Gaussian elimination
and the Schur-complement form of Sylvester's criterion.

\section{Toeplitz matrices: exact symbol characterization}

\begin{theorem}[Bilateral and one-sided Toeplitz coercivity]\label{thm:toeplitz}
Assume \eqref{eq:wiener}, and let $a_J$ be defined by
\eqref{eq:toeplitzform}, where $J=\mathbb Z$ or $J=\N$.  Put
\[
 g_*:=\min_{\theta\in[0,2\pi]}
       \inf_{\|x\|_{\K}=1}\langle G(\theta)x,x\rangle.
\]
Then
\begin{equation}\label{eq:optimal}
 c(a_J)=g_*.
\end{equation}
In particular, the following are equivalent:
\begin{enumerate}[label=\textup{(\roman*)}]
\item $a_J$ is coercive;
\item there is $\gamma>0$ such that $G(\theta)\ge\gamma I_{\K}$ for every
      $\theta\in[0,2\pi]$.
\end{enumerate}
The same optimal constant works for the bilateral and one-sided forms.
\end{theorem}

\begin{proof}
Norm-uniform convergence in \eqref{eq:symbols} makes $F$ and $G$ continuous.
The function
\[
 h(\theta):=\inf_{\|x\|=1}\langle G(\theta)x,x\rangle
\]
is continuous because
$|h(\theta)-h(\varphi)|\le\|G(\theta)-G(\varphi)\|$; hence $g_*$ is a true
minimum.

First consider finitely supported $u,v$ and extend a one-sided vector by zero
to $\mathbb Z$.  With
\[
 (\mathcal Fu)(\theta):=\sum_{n\in\mathbb Z}u_ne^{in\theta},
\]
orthogonality of exponentials and the convention $B_{m-n}$ give
\begin{equation}\label{eq:fourieridentity}
 a_J(u,v)=\frac1{2\pi}\int_0^{2\pi}
 \langle F(\theta)(\mathcal Fu)(\theta),
                (\mathcal Fv)(\theta)\rangle\,d\theta.
\end{equation}
The absolute coefficient sum and vector-valued Young inequality show that the
convolution operator is bounded on $\ell^2(\mathbb Z;\K)$; the one-sided
operator is a compression.  Thus \eqref{eq:fourieridentity} extends to every
$u,v\in\ell^2(J;\K)$.

Taking real parts and using Parseval,
\begin{align*}
 \Real a_J(u,u)
 &=\frac1{2\pi}\int_0^{2\pi}
   \langle G(\theta)(\mathcal Fu)(\theta),
                  (\mathcal Fu)(\theta)\rangle\,d\theta\\
 &\ge g_*\frac1{2\pi}\int_0^{2\pi}
                  \|(\mathcal Fu)(\theta)\|^2\,d\theta
 =g_*\|u\|^2.
\end{align*}
Therefore $c(a_J)\ge g_*$.

For the reverse inequality fix $\theta_0\in[0,2\pi]$, a unit vector
$x\in\K$, and $0<r<1$.  In either index set use the vector supported on
$\N$ given by
\[
 u_n^{(r)}=\sqrt{1-r^2}\,r^n e^{-in\theta_0}x,
 \qquad n\in\N.
\]
It has norm one, and
\[
 \|\mathcal Fu^{(r)}(\theta)\|^2
 =\frac{1-r^2}{|1-re^{i(\theta-\theta_0)}|^2}
 =:P_r(\theta-\theta_0),
\]
the Poisson kernel.  Formula \eqref{eq:fourieridentity} gives
\[
 \Real a_J(u^{(r)},u^{(r)})
 =\frac1{2\pi}\int_0^{2\pi}P_r(\theta-\theta_0)
       \langle G(\theta)x,x\rangle\,d\theta.
\]
The scalar function under the Poisson average is continuous, so the
approximate-identity property yields, as $r\uparrow1$,
\[
 \Real a_J(u^{(r)},u^{(r)})
 \longrightarrow\langle G(\theta_0)x,x\rangle.
\]
Consequently
$c(a_J)\le\langle G(\theta_0)x,x\rangle$ for every $\theta_0$ and every unit
$x$.  Taking the two infima gives $c(a_J)\le g_*$.  Combined with the first
inequality this proves \eqref{eq:optimal} and both equivalences.
\end{proof}

\begin{remark}[Why the one-sided case is not weaker]
For a compression, pointwise symbol positivity is automatically sufficient,
but necessity can fail for arbitrary compression problems.  Here the Hardy
geometric vectors concentrate at any chosen boundary point.  They recover
every pointwise Rayleigh quotient and therefore prevent an artificial gain in
the coercivity constant.
\end{remark}

\section{Incidental correction to the source}

The open question is preceded by Proposition 17 in the source.  That
proposition is not needed above and is false as written.

\begin{proposition}\label{prop:counterexample}
Condition (1.22) of source Proposition 17 is not necessary for coercivity,
even when both component spaces are one-dimensional.
\end{proposition}

\begin{proof}
Take $V_1=V_2=\mathbb C$,
\[
 a_{11}(z,w)=z\overline w,\qquad
 a_{22}(z,w)=z\overline w,\qquad a_{12}=a_{21}=0.
\]
Both optimal diagonal coercivity constants are $\alpha_1=\alpha_2=1$.  The
full form satisfies
\[
a((z_1,z_2),(z_1,z_2))=|z_1|^2+|z_2|^2,
\]
so its optimal coercivity constant is one.  If source condition (1.22) held
for some $\alpha>0$, setting
$z_1=1$ and $z_2=0$ would give
\[
 0\ge\alpha(|z_1|^2+|z_2|^2)-2|z_1||z_2|=\alpha,
\]
a contradiction.
\end{proof}

The error is structural: independently rescaling the two coordinates changes
the norm at which coercivity is tested.  Optimal lower bounds for the diagonal
forms do not convert that family of rescalings into an equivalence.  A separate
proof-gap packet records the exact source statement and counterexample.

\section{Verification and computational checks}

The proof was adversarially checked at the following points:
\begin{enumerate}[leftmargin=*]
\item Multiplying the three factors in \eqref{eq:ldl} recovers the block
matrix with lower-right block $Q$.
\item The Fourier convention in \eqref{eq:fourieridentity} selects
$k+n=m$, hence exactly the coefficient $B_{m-n}$.
\item The normalized geometric vector has norm one and squared Fourier norm
$P_r$; it is valid in both index sets because it is supported in $\N$.
\item Infinite-dimensional $\K$ causes no compactness issue: the infimum over
unit vectors need not be attained, and the proof only takes an infimum.  The
outer minimum exists by operator-norm continuity.
\end{enumerate}

The diagnostic script \texttt{code/check\_coercivity.py} checks 100 random
finite Schur factorizations, a $2\times2$ matrix-valued banded Toeplitz symbol
against a 384-block one-sided compression, and the source counterexample.  It
reports
\[
 g_{\rm grid}=2.475765077,\qquad
 \lambda_{\min}(T_{384}^{\rm herm})=2.475784964,
\]
consistent with finite compressions approaching the symbol edge from above.
These computations are diagnostics only; no numerical statement is used in
the proof.

\section{Novelty search, limitations, and review status}

A bounded check on 9 August 2026 searched the run indexes and arXiv-focused
web results using the exact source wording and the phrases
\emph{operator-valued Toeplitz coercivity symbol},
\emph{Toeplitz sesquilinear forms coercive}, and
\emph{block operator matrix Schur complement coercivity}.  No later paper
explicitly claiming to resolve Cardanobile's question was found.  General
Toeplitz positivity and Schur-complement theory are classical; see, for
example, \cite{BottcherSilbermann}.  Accordingly, this packet claims a new
self-contained identification and proof relative to the source question, not
priority for the ingredients or a comprehensive literature novelty claim.

The Toeplitz theorem assumes the natural Wiener condition
\eqref{eq:wiener}.  It therefore covers every finite-band form matrix and a
large norm-continuous class, but does not claim to classify arbitrary weakly
specified Toeplitz arrays.  Since the source asks whether such translations
are possible and names the two structures rather than demanding maximal
generality, the packet classifies the result as a full positive answer.  A
reviewer adopting the stronger historical interpretation may conservatively
reclassify it as a substantial partial result.

\medskip
\noindent\textbf{Human-review recommendation.}
Send to human review.  The main proof checks are the Toeplitz indexing, the
Hardy/Poisson necessity argument, and the intended scope of the informal
source question.  The mathematical theorem itself has no conditional
dependency.

\begin{thebibliography}{9}
\bibitem{Cardanobile}
S. Cardanobile,
\emph{Diffusion systems and heat equations on networks},
PhD thesis, arXiv:0807.2362, 2008.

\bibitem{BottcherSilbermann}
A. B\"ottcher and B. Silbermann,
\emph{Analysis of Toeplitz Operators},
2nd ed., Springer Monographs in Mathematics, Springer, 2006.
\end{thebibliography}

\end{document}
